\section{Implementation and independent verification}
\label{sec:implementation}

\subsection{Module structure}

The six new runtime modules are listed below. Their only dependencies are
Python's standard library and the pinned existing \code{fastunknot}
modules. Regina is optional and used by the external audit, not by the
new producer or its certificate checker.

\begin{center}\small
\begin{tabular}{>{\raggedright\arraybackslash}p{0.45\linewidth}>{\raggedright\arraybackslash}p{0.47\linewidth}}
\toprule
Module under \path{fastunknot/} & Responsibility\\
\midrule
\path{orbit_transversal.py} & Forward original-label interval tracking and
least-transversal production.\\
\path{orbit_transversal_verify.py} & Independent reverse gap-insertion
reconstruction after source-bound trace replay.\\
\path{topology_spectrum.py} & Sparse cover inversion, surface classification,
normal scaling, and independent forward-equation checkers.\\
\path{normal_topology_geometry.py} & Boundary inclusion, integral Euler and
marker weights, vertex-link classification.\\
\path{normal_topology.py} & Validated three-query producer, optional
quadrilateral-core reduction, and command-line entry point.\\
\path{normal_topology_verify.py} & Original-source core reconstruction,
independent local replays, spectrum and restoration verification.\\
\bottomrule
\end{tabular}
\end{center}

\Needspace{18\baselineskip}
\subsection{The public API}

The main call is
\begin{lstlisting}[language=Python]
from fastunknot.normal_topology import normal_topology_spectrum
from fastunknot.normal_topology_verify import verify_normal_topology_spectrum

answer = normal_topology_spectrum(
    triangulation, coordinates,
    reduce_core=True,
    max_cycles=None,
    periodic_rule="fine_wilf",
    record_certificate=True,
)
assert answer["status"] == "COMPLETE"
assert verify_normal_topology_spectrum(
    triangulation, coordinates, answer["certificate"]
)
\end{lstlisting}
The output has \code{topology\_spectrum} rows with fields \code{chi},
\code{boundary\_components}, \code{orientable}, \code{multiplicity},
and exactly one of \code{genus} and \code{crosscaps}. It also gives total
components, orientable and nonorientable counts, boundary circles and Euler
characteristic. These mathematical claims concern the original source,
even when the statistics describe smaller core queries.

The public \code{reduce\_core=False} setting exposes the direct observer for
comparison and for inputs on which preprocessing is not helpful.
The two inherited periodic rules are available. The default uses the
existing Fine--Wilf merger; \code{aht} uses the older threshold. Both
represent the same orbit relation, and the resulting least transversal and
type spectrum must agree.

The certificate schema is \code{normal-topology-spectrum-v1}. It binds the
exact normalized coordinate vector and face records by a digest, retains
the requested reduction mode, and includes the divisor, peeled vertex-link
multiplicities, exact core vector, boundary-transversal proof, two weighted
proofs, the core spectrum, and the restored original summary. The pure
vertex-link case has divisor zero and no orbit-query proof: validated link
classification and exact coordinate decomposition suffice.

\subsection{Verification architecture and trusted boundary}

The geometry validator, normal-arc conventions, cell ownership, and exact
integer codec are shared trusted reconstruction code. Producer/checker
independence does not mean that two unrelated triangulation engines are
embedded in the runtime. Independent checking is instead structured around
the vulnerable arithmetic and reduction layers:
\begin{enumerate}
\item The original triangulation and coordinates are revalidated, and the
original source digest is checked.
\item The checker rebuilds corner groups, independently recomputes minima
and quadrilateral gcd, and checks exact divisibility and all supplied core
fields. It does not call the producer's core routine.
\item The boundary orbit proof is replayed against its actual reconstructed
source. The checker reconstructs its transversal backwards rather than
calling the forward producer.
\item Both weighted proofs are checked by the existing independent weight
transport checker against freshly reconstructed two-coordinate weights.
\item The claimed core spectrum is checked by \emph{forward} reconstruction
of the entire base and cover histograms, without invoking sparse inversion.
\item The restored spectrum is checked with independent typewise scaling
equations, and all original totals are reconstructed independently.
\end{enumerate}
Tests disable producer discovery, forward selector construction, producer
weight replay, spectrum inversion, and producer scaling while verifying a
completed certificate. Separate Regina comparisons provide an external
check on the shared geometry boundary for bounded cases.

A source digest alone is not treated as a mathematical proof. It prevents
accidental reassociation of a proof with another input, while each local
certificate is still checked against the actual reconstructed relation and
weights. Mutating the original peeled vertex-link multiplicity while keeping
the same core must fail verification.

\subsection{Resource and transport contracts}

One shared \code{max\_cycles} budget covers boundary, full-surface, and
doubled-surface orbit discoveries. If a call exhausts this allowance, it
returns \code{INCONCLUSIVE} with stage and statistics, and no completed
spectrum, component count, or certificate. An allowance of exactly the
sum of the required cycles succeeds. The cycle budget is not silently
reinterpreted as a bound on validation or weight transport.

A cooperative \code{check()} callback is polled throughout validation,
discovery, interval compilation, weight transport, algebraic recovery and
verification. Its exceptions propagate. False-valued callable objects are
respected rather than replaced by a default callback. The command-line
wrapper exposes a wall-clock deadline through this mechanism and reports
an inconclusive result on timeout.

Arbitrary-size integers remain supported without changing Python's global
decimal-conversion limit. The existing \code{json\_safe} codec serializes
large integers as signed hexadecimal strings; certificate checkers decode
the declared integer fields exactly and reject Boolean values in their
place. Large output multiplicities remain a single encoded integer.

An arbitrary caller-supplied valid trace need not have the maintained
producer's short length. Verification and transversal compilation are
polynomial in the explicit certificate size; the short-trace guarantee
belongs to the maintained discovery algorithm. This distinction is necessary
for honest time accounting on adversarial serialized inputs.

\subsection{Command-line use and integration}

From the delivered \path{code/fast} directory, an input JSON file contains
\code{triangulation} and \code{coordinates}. The two commands are
\begin{lstlisting}
python3 -m fastunknot.normal_topology census \
  topology_research/data/klein_torus.json --certificate -o klein_result.json

python3 -m fastunknot.normal_topology verify \
  topology_research/data/klein_torus.json klein_result.json
\end{lstlisting}
Add \code{--direct} to disable the core, or \code{--max-cycles} and
\code{--timeout} to bound a discovery run. Verification accepts either a
certificate object or a complete result containing one.

The repository integration patch adds the runtime modules, their focused
tests, and the reproducible research inputs/drivers. It is prepared against
the pinned baseline; no remote commit or pull request is created by this
delivery. The complete existing recognition pipeline is not wired to invent
or accept normal vectors automatically. The new API is a verified observer
that future surface-discovery and hierarchy code can call explicitly.
